% TOP_PROBLEM: {"id": 4700003, "title": "Systems with homogeneous components II", "category_id": 11, "category": {"id": 11, "name": "dynamical_systems", "display_name": "Dynamical Systems", "description": "Problems about long-term behavior of deterministic systems, Hamiltonian dynamics, and periodic orbits.", "slug": "dynamical-systems", "order_index": 11, "created_at": "2026-07-31T15:26:25.671Z"}, "status": "open", "problem_number": "AMR-046-0003", "set_id": 13}
% FIRST_POSTED: 2026-10-02
% ENTRY_KIND: statement_only
% UnsolvedMath ID 4700003; source catalogue code AMR-046-0003
% !TeX program = lualatex
% Definition and sources only; this file is not an LLM solution attempt.
\documentclass[11pt,a4paper]{article}
\usepackage[margin=25mm]{geometry}
\usepackage{fontspec}
\setmainfont{lmroman10-regular.otf}[BoldFont=lmroman10-bold.otf,
 ItalicFont=lmroman10-italic.otf,BoldItalicFont=lmroman10-bolditalic.otf]
\usepackage{amsmath,amssymb,mathtools}
\usepackage{unicode-math}
\setmathfont{latinmodern-math.otf}
\AtBeginDocument{\let\setminus\smallsetminus}
\usepackage{xurl,hyperref}
\hypersetup{colorlinks=true,urlcolor=blue}
\setcounter{secnumdepth}{0}
\setlength{\parindent}{0pt}
\setlength{\parskip}{5pt}
\setlength{\emergencystretch}{3em}
\begin{document}
\section*{Systems with homogeneous components II}\label{4700003}
{\small UnsolvedMath ID 4700003 · AMR-046-0003 · Dynamical Systems · AMR Open Problem Lists}
\subsection{Definitions and mathematical statement}
Let $a,b,c,d,e,f\in\mathbb R$ and consider
\[
 \dot x=ax+by,\qquad
 \dot y=cx^3+dx^2y+exy^2+fy^3.
\]
A limit cycle is an isolated nonconstant periodic orbit of this autonomous
system. Different starting points on one closed orbit represent the same
cycle. The first question asks whether the supremum of the number of such
cycles, over all parameter choices, equals two. This is a global count in
$\mathbb R^2$, including cycles not obtainable from small perturbations of
an equilibrium.

The source singles out the one-parameter subfamily
\[
 \dot x=y,\qquad \dot y=-x^3+dx^2y+y^3,
 \qquad d\in\mathbb R.
\]
Its second question asks for a simple proof that this system has at most one
limit cycle for each $d$, if that assertion is true. Here ``uniqueness'' is
conditional on existence; it does not assert that every real parameter yields
a periodic orbit. A proof should cover the entire parameter line, including
the intervals between the partial existence and nonexistence ranges in the
source.

Both questions concern isolation among periodic orbits, so a center with a
continuum of closed trajectories does not supply infinitely many cycles.
The special subfamily offers a more concrete first test of the unequal-degree
homogeneous-component problem, but resolving it alone does not prove the
two-cycle bound for the full six-parameter family.
\subsection{Short English statement}
Determine whether the full linear–cubic family has at most two limit cycles, and whether the displayed subfamily has at most one.
\subsection{Sources}
\begin{itemize}
\item[S1] ULAM.AI and the UnsolvedMath Contributors, supplied problems.json, record 4700003 (AMR-046-0003), AMR Open Problem Lists. Catalogue identity and source transcription; CC BY 4.0.
\url{https://huggingface.co/datasets/ulamai/UnsolvedMath}.
\item[S2] Gasull - Some open problems in low dimensional dynamical systems (2020), Problem environment 3: Systems with homogeneous components II. Original source indexed by AMR-046-0003.
\url{https://arxiv.org/abs/2012.02524}.
\end{itemize}
\end{document}
